The second stage performs switch allocation. Only input VCs with valid first-stage information are allowed to request access to the switch. Since multiple input VCs may request the same output port in the same cycle, arbitration is required.

The switch allocator uses the separable input-first allocation method described earlier. For each input port, at most one local VC can win the input-side arbitration. For each output port, at most one requester can win the output-side arbitration. When a flit wins switch allocation, it is dequeued from the selected input VC FIFO and stored in the second-stage pipeline registers together with its control information, selected output port, and output VC.

When the flit moves from this stage to the next stage, one downstream credit is consumed. If the dequeued flit is a tail or head-tail flit, the corresponding input-VC lock is cleared at the end of this stage, allowing the next packet in the same input VC to perform routing computation and virtual-channel allocation later.